
Reproducibility assessment tools operate after experiments run, providing no guidance for predicting which datasets will yield reproducible results before writing code. This work investigates whether dataset metadata completeness predicts reproducibility variance in machine learning benchmarks. Analyzing 300 datasets with 21,312 matched experimental runs, results indicate that top-quartile metadata completeness is associated with a 42.1\% reduction in interquartile range (IQR) of performance outcomes compared to bottom-quartile datasets (95\% CI: 39.1--51.7\%). Mediation analysis identifies preprocessing entropy as a candidate mechanism, accounting for 64.7\% of the observed association (Sobel Z = 16.02, p < 0.0001): datasets with more complete documentation exhibit lower Shannon entropy in preprocessing component distributions, which corresponds to reduced performance variance. The association persists in early-run subsamples (90.7\% preservation) and within single algorithm families (p < 0.001), addressing concerns about reverse causality and algorithm-mix confounding. All experiments were conducted on synthetic data due to API availability constraints; validation with real-world data remains necessary. Code and data available at [repository].

